%! Unit Opener: Why Algebra Fluency Matters — Unit 1, Lesson 1.0 — Slides (Beamer)
\documentclass[aspectratio=169, 11pt]{beamer}
\usepackage{atda-beamer}   % provides \royalheader{} and \sectionlabel[color]{}

\begin{document}

% TITLE SLIDE (CourseName is not defined in beamer, so the name is literal)
{
\setbeamercolor{background canvas}{bg=royal}
\begin{frame}[plain]
  \vspace{2.8cm}\hspace{0.55cm}%
  \begin{minipage}{0.9\textwidth}
    {\color{goldacc}\bfseries\small LESSON 1.0}\\[0.5cm]
    {\color{white}\bfseries\LARGE Unit Opener: Why Algebra Fluency Matters}\\[1.2cm]
    {\color{mistmid}\small Algebra, Trigonometry, and Data Analysis~$\cdot$~Unit 1}
  \end{minipage}
\end{frame}
}

% ── HOOK ──────────────────────────────────────────────────────────────────────
\begin{frame}[t, plain]
  \royalheader{The Fourth of July Shell}
  \sectionlabel[goldacc]{hook}
  A shell is launched from a platform on a $64$-foot cliff. Its height above the water,
  in feet, $t$ seconds after launch:
  \[ h(t) = -16t^2+48t+64 \qquad\text{and also}\qquad h(t) = -16(t-4)(t+1) \]
  \begin{block}{These are the same function.}
    The \textbf{safety officer} needs to know when the shell hits the water.\\
    The \textbf{announcer} needs to know how high the platform is.\\[2pt]
    Which form would you hand to each one --- and why?
  \end{block}
\end{frame}

% ── ARE THEY REALLY THE SAME? ─────────────────────────────────────────────────
\begin{frame}[t, plain]
  \royalheader{First: Are They Really the Same?}
  \sectionlabel{verify by expanding}
  \[
  \begin{aligned}
    -16(t-4)(t+1) &= -16\big(t^2+t-4t-4\big) \\
                  &= -16\big(t^2-3t-4\big) \\
                  &= -16t^2+48t+64
  \end{aligned}
  \]
  \begin{block}{The error to avoid}
    The $-16$ reaches \textbf{every} term inside the parentheses --- not just the first one.
    This expand-and-compare check verifies every factoring answer you will write this year.
  \end{block}
\end{frame}

% ── TWO FORMS, TWO QUESTIONS ──────────────────────────────────────────────────
\begin{frame}[t, plain]
  \royalheader{Two Forms, Two Questions}
  \sectionlabel{what each form shows}
  \renewcommand{\arraystretch}{1.5}
  \begin{tabularx}{\textwidth}{@{}p{4.6cm} p{2.6cm} X@{}}
    \textbf{The question} & \textbf{Form} & \textbf{What you read off} \\
    \hline
    How high is the platform? & expanded & the constant, $64$ feet \\
    When does it hit the water? & factored & the zero, $t=4$ seconds \\
  \end{tabularx}
  \vspace{0.5cm}

  Neither form is ``better.'' \textbf{Fluency is being able to move between them on demand.}
\end{frame}

% ── THE ZERO THAT DOESN'T COUNT ───────────────────────────────────────────────
\begin{frame}[t, plain]
  \royalheader{Two Zeros, One Answer}
  \sectionlabel{interpret in context}
  The factored form gives two zeros: $\;t=4\;$ and $\;t=-1$.
  \vspace{0.3cm}
  \begin{itemize}
    \setlength{\itemsep}{0.35cm}
    \item Both are real numbers. Both make $h(t)=0$.
    \item Only $t=4$ describes the shell --- $t=-1$ is one second \emph{before} the launch.
    \item Every model has a domain the situation allows. Compute, then ask:
          \textbf{does this answer mean anything here?}
  \end{itemize}
\end{frame}

% ── ROADMAP ───────────────────────────────────────────────────────────────────
\begin{frame}[t, plain]
  \royalheader{The Unit 1 Roadmap}
  \sectionlabel{where we are going}
  \renewcommand{\arraystretch}{1.25}
  \small
  \begin{tabularx}{\textwidth}{@{}c p{2.7cm} X@{}}
    \textbf{Lesson} & \textbf{Standard} & \textbf{What you will be able to do} \\
    \hline
    1.1 & \texttt{A2.EO.3a}    & Exponent rules; add, subtract, multiply polynomials \\
    1.2 & \texttt{A2.EO.2a--c} & Radicals and rational exponents \\
    1.3 & \texttt{A2.EO.3b}    & Factor with a GCF and by grouping \\
    1.4 & \texttt{A2.EO.3b}    & Factor trinomials \\
    1.5 & \texttt{A2.EO.3b, d} & Special forms and polynomial identities \\
    1.6 & \texttt{A2.EI.2b, EI.6a--b} & Solve polynomial equations by factoring \\
    1.7 & \texttt{A2.EO.1a--b} & Simplify rational expressions \\
  \end{tabularx}
\end{frame}

% ── WHERE THE ALGEBRA GETS SPENT ──────────────────────────────────────────────
\begin{frame}[t, plain]
  \royalheader{Where This Algebra Gets Spent}
  \sectionlabel[goldacc]{the bridge}
  \begin{itemize}
    \setlength{\itemsep}{0.4cm}
    \item \textbf{Exponent rules} (1.1) become the \textbf{properties of logarithms} in Unit 3:
          $\log(ab)=\log a+\log b$ is the same idea written backwards.
    \item \textbf{Factoring} (1.3--1.5) is what makes equations solvable in 1.6 --- and every
          trigonometric equation in Unit 8 ends the same way.
    \item \textbf{Simplifying quotients} (1.7) reduces the counting and probability formulas of
          Unit 4 to numbers you can compute by hand.
  \end{itemize}
\end{frame}

% ── ACTIVITY LAUNCH ───────────────────────────────────────────────────────────
\begin{frame}[t, plain]
  \royalheader{Group Activity: The Spirit-Wear Sale}
  \sectionlabel{groups of three --- everyone starts at Tier R}
  At a price of $x$ dollars per shirt, the robotics club expects to sell $120-4x$ shirts.
  Revenue, in dollars:
  \[ R(x) = x(120-4x) \]
  \begin{block}{Your job}
    Expand it. Find its zeros. Then explain, in context, why the club takes in \emph{nothing} at
    \textbf{both} of those prices --- for two different reasons.
  \end{block}
\end{frame}

% ── CLOSE ─────────────────────────────────────────────────────────────────────
\begin{frame}[t, plain]
  \royalheader{Exit Ticket}
  \sectionlabel[goldacc]{notes closed --- 5 minutes}
  \begin{enumerate}
    \setlength{\itemsep}{0.35cm}
    \item Expand $3x(x-4)$ and compare it with $3x^2-12x$.
    \item Read the zeros off $H(t)=-2(t-9)(t+1)$ and circle the one that describes the flight.
    \item ``The factored form is just the better form.'' Do you agree? Name a question the
          \textbf{expanded} form answers faster.
  \end{enumerate}
  \vspace{0.3cm}
  {\color{royal}\bfseries Next: Lesson 1.1 --- Simplifying Expressions and Exponent Rules.}
\end{frame}

\end{document}
